\chapter*{LAMPIRAN}

\section{Skema Lengkap \emph{Knowledge Graph}}
\label{app:schema}

Tabel~\ref{tab:full-schema} mencantumkan setiap tripel label-relasi-label yang terdapat
pada basis data \texttt{sihdfhir} yang dimuat, sebagaimana dilaporkan oleh penyampelan skema
basis data itu sendiri. Label struktural \texttt{:Resource} dihilangkan di seluruh tabel:
label tersebut melekat pada setiap simpul kecuali \texttt{:Code} sehingga menyertakannya akan
menambahkan duplikat bagi hampir setiap baris tanpa memberikan informasi apa pun (lihat
Subbab~\ref{sec:kg-construction}). Tersisa 24 label dan 95 tripel. Tripel yang label dan tipe
relasinya ditelusuri oleh kueri rekomendasi---yaitu subgraf \emph{terpakai} pada
Gambar~\ref{fig:graph-model}---merupakan subhimpunan dari tabel ini.

\begin{longtable}{@{}p{46mm}p{52mm}p{40mm}@{}}
\caption{Skema lengkap \emph{knowledge graph} yang dimuat (\texttt{:Resource} dihilangkan).
Setiap isian pada kolom \emph{Dari} dan \emph{Ke} merupakan label simpul dan setiap isian pada
kolom \emph{Relasi} merupakan tipe relasi sehingga label di sini ditulis tanpa titik dua
pendahulu yang menandainya di dalam teks.}
\label{tab:full-schema} \\
\hline
\textbf{Dari} & \textbf{Relasi} & \textbf{Ke} \\
\hline
\endfirsthead
\multicolumn{3}{l}{\small\itshape Tabel \thetable\ lanjutan dari halaman sebelumnya} \\
\hline
\textbf{Dari} & \textbf{Relasi} & \textbf{Ke} \\
\hline
\endhead
\hline
\multicolumn{3}{r}{\small\itshape bersambung ke halaman berikutnya} \\
\endfoot
\hline
\endlastfoot
\texttt{Abstract} & \texttt{MENTIONS} & \texttt{Concept} \\
\texttt{AllergyIntolerance} & \texttt{CLINICAL\_STATUS} & \texttt{Code} \\
\texttt{AllergyIntolerance} & \texttt{CODE} & \texttt{Code} \\
\texttt{AllergyIntolerance} & \texttt{PATIENT} & \texttt{Patient} \\
\texttt{AllergyIntolerance} & \texttt{REACTION\_MANIFESTATION} & \texttt{Code} \\
\texttt{AllergyIntolerance} & \texttt{VERIFICATION\_STATUS} & \texttt{Code} \\
\texttt{CarePlan} & \texttt{ACTIVITY\_DETAIL\_CODE} & \texttt{Code} \\
\texttt{CarePlan} & \texttt{ACTIVITY\_DETAIL\_REASON} & \texttt{Condition} \\
\texttt{CarePlan} & \texttt{ACTIVITY\_DETAIL\_REASON\_CODE} & \texttt{Code} \\
\texttt{CarePlan} & \texttt{ADDRESSES} & \texttt{Condition} \\
\texttt{CarePlan} & \texttt{CARE\_TEAM} & \texttt{CareTeam} \\
\texttt{CarePlan} & \texttt{CATEGORY} & \texttt{Code} \\
\texttt{CarePlan} & \texttt{ENCOUNTER} & \texttt{Encounter} \\
\texttt{CarePlan} & \texttt{SUBJECT} & \texttt{Patient} \\
\texttt{CareTeam} & \texttt{ENCOUNTER} & \texttt{Encounter} \\
\texttt{CareTeam} & \texttt{MANAGING\_ORGANIZATION} & \texttt{Organization} \\
\texttt{CareTeam} & \texttt{PARTICIPANT\_MEMBER} & \texttt{Organization} \\
\texttt{CareTeam} & \texttt{PARTICIPANT\_MEMBER} & \texttt{Patient} \\
\texttt{CareTeam} & \texttt{PARTICIPANT\_MEMBER} & \texttt{Practitioner} \\
\texttt{CareTeam} & \texttt{PARTICIPANT\_ROLE} & \texttt{Code} \\
\texttt{CareTeam} & \texttt{REASON\_CODE} & \texttt{Code} \\
\texttt{CareTeam} & \texttt{SUBJECT} & \texttt{Patient} \\
\texttt{Code} & \texttt{CO\_OCCURS\_IN\_LITERATURE} & \texttt{Code} \\
\texttt{Component} & \texttt{CODE} & \texttt{Code} \\
\texttt{Concept} & \texttt{SAME\_CONCEPT} & \texttt{Code} \\
\texttt{Condition} & \texttt{CATEGORY} & \texttt{Code} \\
\texttt{Condition} & \texttt{CLINICAL\_STATUS} & \texttt{Code} \\
\texttt{Condition} & \texttt{CODE} & \texttt{Code} \\
\texttt{Condition} & \texttt{ENCOUNTER} & \texttt{Encounter} \\
\texttt{Condition} & \texttt{SUBJECT} & \texttt{Patient} \\
\texttt{Condition} & \texttt{VERIFICATION\_STATUS} & \texttt{Code} \\
\texttt{Device} & \texttt{PATIENT} & \texttt{Patient} \\
\texttt{Device} & \texttt{TYPE} & \texttt{Code} \\
\texttt{Encounter} & \texttt{CLASS} & \texttt{Code} \\
\texttt{Encounter} & \texttt{HOSPITALIZATION\_DISCHARGE\_DISPOSITION} & \texttt{Code} \\
\texttt{Encounter} & \texttt{LOCATION} & \texttt{Location} \\
\texttt{Encounter} & \texttt{PARTICIPANT\_INDIVIDUAL} & \texttt{Practitioner} \\
\texttt{Encounter} & \texttt{PARTICIPANT\_TYPE} & \texttt{Code} \\
\texttt{Encounter} & \texttt{REASON\_CODE} & \texttt{Code} \\
\texttt{Encounter} & \texttt{SERVICE\_PROVIDER} & \texttt{Organization} \\
\texttt{Encounter} & \texttt{SUBJECT} & \texttt{Patient} \\
\texttt{Encounter} & \texttt{TYPE} & \texttt{Code} \\
\texttt{Guideline} & \texttt{CONTAINS} & \texttt{Recommendation} \\
\texttt{ImagingStudy} & \texttt{ENCOUNTER} & \texttt{Encounter} \\
\texttt{ImagingStudy} & \texttt{LOCATION} & \texttt{Location} \\
\texttt{ImagingStudy} & \texttt{PROCEDURE\_CODE} & \texttt{Code} \\
\texttt{ImagingStudy} & \texttt{SERIES\_BODY\_SITE} & \texttt{Code} \\
\texttt{ImagingStudy} & \texttt{SERIES\_INSTANCE\_SOP\_CLASS} & \texttt{Code} \\
\texttt{ImagingStudy} & \texttt{SERIES\_MODALITY} & \texttt{Code} \\
\texttt{ImagingStudy} & \texttt{SUBJECT} & \texttt{Patient} \\
\texttt{Immunization} & \texttt{ENCOUNTER} & \texttt{Encounter} \\
\texttt{Immunization} & \texttt{LOCATION} & \texttt{Location} \\
\texttt{Immunization} & \texttt{PATIENT} & \texttt{Patient} \\
\texttt{Immunization} & \texttt{VACCINE\_CODE} & \texttt{Code} \\
\texttt{Location} & \texttt{PHYSICAL\_TYPE} & \texttt{Code} \\
\texttt{Medication} & \texttt{CODE} & \texttt{Code} \\
\texttt{MedicationAdministration} & \texttt{CONTEXT} & \texttt{Encounter} \\
\texttt{MedicationAdministration} & \texttt{MEDICATION} & \texttt{Code} \\
\texttt{MedicationAdministration} & \texttt{MEDICATION} & \texttt{Medication} \\
\texttt{MedicationAdministration} & \texttt{REASON} & \texttt{Condition} \\
\texttt{MedicationAdministration} & \texttt{REASON\_CODE} & \texttt{Code} \\
\texttt{MedicationAdministration} & \texttt{SUBJECT} & \texttt{Patient} \\
\texttt{MedicationRequest} & \texttt{CATEGORY} & \texttt{Code} \\
\texttt{MedicationRequest} & \texttt{DOSAGE\_INSTRUCTION\_ADDITIONAL\_INSTRUCTION} & \texttt{Code} \\
\texttt{MedicationRequest} & \texttt{DOSAGE\_INSTRUCTION\_DOSE\_AND\_RATE\_TYPE} & \texttt{Code} \\
\texttt{MedicationRequest} & \texttt{ENCOUNTER} & \texttt{Encounter} \\
\texttt{MedicationRequest} & \texttt{MEDICATION} & \texttt{Code} \\
\texttt{MedicationRequest} & \texttt{MEDICATION} & \texttt{Medication} \\
\texttt{MedicationRequest} & \texttt{REASON} & \texttt{Condition} \\
\texttt{MedicationRequest} & \texttt{REASON\_CODE} & \texttt{Code} \\
\texttt{MedicationRequest} & \texttt{REQUESTER} & \texttt{Practitioner} \\
\texttt{MedicationRequest} & \texttt{SUBJECT} & \texttt{Patient} \\
\texttt{Observation} & \texttt{CATEGORY} & \texttt{Code} \\
\texttt{Observation} & \texttt{CODE} & \texttt{Code} \\
\texttt{Observation} & \texttt{COMPONENT} & \texttt{Component} \\
\texttt{Observation} & \texttt{ENCOUNTER} & \texttt{Encounter} \\
\texttt{Observation} & \texttt{SUBJECT} & \texttt{Patient} \\
\texttt{Observation} & \texttt{VALUE} & \texttt{Code} \\
\texttt{Organization} & \texttt{TYPE} & \texttt{Code} \\
\texttt{Patient} & \texttt{COMMUNICATION\_LANGUAGE} & \texttt{Code} \\
\texttt{Patient} & \texttt{EXTENSION\_VALUE} & \texttt{Code} \\
\texttt{Patient} & \texttt{IDENTIFIER\_TYPE} & \texttt{Code} \\
\texttt{Patient} & \texttt{MARITAL\_STATUS} & \texttt{Code} \\
\texttt{PractitionerRole} & \texttt{CODE} & \texttt{Code} \\
\texttt{PractitionerRole} & \texttt{SPECIALTY} & \texttt{Code} \\
\texttt{Procedure} & \texttt{CODE} & \texttt{Code} \\
\texttt{Procedure} & \texttt{ENCOUNTER} & \texttt{Encounter} \\
\texttt{Procedure} & \texttt{LOCATION} & \texttt{Location} \\
\texttt{Procedure} & \texttt{REASON} & \texttt{Condition} \\
\texttt{Procedure} & \texttt{REASON\_CODE} & \texttt{Code} \\
\texttt{Procedure} & \texttt{SUBJECT} & \texttt{Patient} \\
\texttt{Recommendation} & \texttt{CITES} & \texttt{Abstract} \\
\texttt{Recommendation} & \texttt{CONTRAINDICATES} & \texttt{Concept} \\
\texttt{Recommendation} & \texttt{RECOMMENDS} & \texttt{Concept} \\
\texttt{Recommendation} & \texttt{TARGETS} & \texttt{Concept} \\
\end{longtable}

\section{Kueri Rekomendasi}
\label{app:queries}
Kueri berikut adalah kueri yang dievaluasi per pasien untuk menghasilkan himpunan
rekomendasi yang berlaku; Listing~\ref{lst:applicable-core} pada Bab~III
memperlihatkan intinya. Kueri ini disalin apa adanya dari modul evaluasi.

\begin{lstlisting}[caption={Rekomendasi yang berlaku bagi satu pasien.},label={lst:applicable-full}]
MATCH (p:Patient {uid:$pid})
WHERE EXISTS { (p)<-[:SUBJECT]-(:Condition)-[:CODE]->(:Code)<-[:SAME_CONCEPT]-(:Concept {umlsCui:$gate}) }
MATCH (rec:Recommendation)
WHERE NOT rec.cor STARTS WITH '3'
OPTIONAL MATCH (rec)-[:TARGETS]->(tc:Concept)
  WHERE tc.umlsCui <> $gate
    AND ( EXISTS { (tc)-[:SAME_CONCEPT]->(:Code)<-[:CODE]-(:Condition) }
       OR EXISTS { (tc)-[:SAME_CONCEPT]->(:Code)<-[:CODE]-(:Procedure) } )
WITH p, rec, collect(DISTINCT tc) AS extraConds
WITH p, rec, extraConds,
     [tc IN extraConds WHERE
        EXISTS { (p)<-[:SUBJECT]-(:Condition)-[:CODE]->(:Code)<-[:SAME_CONCEPT]-(tc) }
        OR EXISTS { (p)<-[:SUBJECT]-(:Procedure)-[:CODE]->(:Code)<-[:SAME_CONCEPT]-(tc) }] AS matched
WHERE size(extraConds) = 0
   OR (coalesce(rec.targetLogic, 'AND') = 'OR'  AND size(matched) >= 1)
   OR (coalesce(rec.targetLogic, 'AND') = 'AND' AND size(matched) = size(extraConds))
RETURN rec.recommendationId AS recId
\end{lstlisting}

\section{Modul dan Keluarannya}
\label{app:modules}
Tabel~\ref{tab:modules} mencantumkan dua belas modul yang membentuk alur utamanya. Beberapa
skrip lain menopang alur tersebut tanpa menjadi bagian darinya: \texttt{cohort.cypher}
mendefinisikan kohort CCD menurut kriteria klinis; \texttt{make\_worksheets.py}
membangkitkan paket pelabelan manual dan \texttt{build\_gold\_labels.py} menggabungkan
lembar kerja yang telah diisi menjadi \texttt{eval/gold\_labels.csv}.

Evaluasi ini bertumpu pada dua standar rujukan dan penamaan berkasnya membedakan keduanya:

\begin{itemize}
  \item \texttt{silver\_oracle.py} mengodekan logika panduan dan melabeli setiap
  pasien dengan membaca \emph{bundle} FHIR mentah, lalu menuliskan
  \texttt{eval/silver\_oracle\_labels.csv}. Keluarannya dibangkitkan, bukan dikurasi.
  \item \texttt{eval/gold\_labels.csv} memuat label manual yang dihasilkan manusia
  dengan membaca teks panduan terhadap rekam medis pasien.
  \item \texttt{crs\_eval.py} menilai sistem terhadap salah satu standar tersebut;
  \texttt{obs\_gate.py} mengevaluasi \emph{observation gate};
  \texttt{pick\_review\_patients.py} memeringkat kohort untuk sampel manual.
  \item \texttt{subjektif\_audit.py} mengukur data subjektif pada rekam medis FHIR ke-30
  pasien sampel---catatan klinis, keluhan utama, gejala jantung, skala fungsional, kuesioner
  yang dilaporkan pasien, dan diagnosis angina---dan menuliskan
  \texttt{eval/subjektif\_audit.txt}, yang menopang keterbatasan tentang data gejala jantung
  pada Subbab~\ref{sec:limitations}.
  \item \texttt{kappa\_bootstrap.py} menaksir interval kepercayaan bagi $\kappa$
  \emph{silver}-\emph{gold} dengan \emph{bootstrap} yang mencuplik ulang pasien maupun sel. Skrip itu juga melaporkan indeks prevalensi, indeks bias, PABAK, dan \emph{jackknife} per pasien;
  keluarannya \texttt{eval/kappa\_bootstrap.txt} menopang Subbab~\ref{sec:kappa-ci}.
  \item \texttt{predicate\_audit.py} memeriksa di mana standar rujukan berbagi predikat
  dengan sistem yang diujinya, membuktikan kesamaan fungsinya atas 200.876 vektor ciri, dan
  melaporkan keselarasan setelah rekomendasi berpredikat bersama disingkirkan; keluarannya
  \texttt{eval/predicate\_audit.txt} menopang Subbab~\ref{sec:shared-predicate}.
  \item \texttt{csv\_gate.py} membangun kembali \emph{observation gate} dari ekspor CSV
  Synthea semata, lalu membandingkannya terhadap jalur \emph{bundle} FHIR dan menuliskan
  \texttt{eval/csv\_gating\_experiment.txt}. Skrip inilah yang menghasilkan angka pada
  Subbab~\ref{sec:csv-gating}.
  \item Direktori \texttt{telaah-klinisi/} memuat perangkat telaah klinisi.
  \texttt{buat\_telaah.py} dan \texttt{perbarui\_lembar.py} membangkitkan lembar telaah
  beserta berkas kuncinya. \texttt{baca\_hasil.py} membaca lembar yang telah diisi menjadi satu
  tabel dan \texttt{analisis\_hasil.py} menghitung proporsi penilaian, interval
  \emph{bootstrap} per pasien, serta kesepakatan antarpenilai yang dilaporkan pada
  Subbab~\ref{sec:clinician}.
\end{itemize}

\subsection{Pemuatan Idempoten}
Listing~\ref{lst:loader-merge} memuat empat pernyataan tulis yang dijalankan
\emph{loader}. Semuanya berupa \texttt{MERGE} dan justru itulah yang memungkinkan sebuah
\emph{bundle} dimuat ulang tanpa menggandakan simpul, kode, maupun sisi---sifat yang
ditegaskan oleh cacah \texttt{:Patient} pada Subbab~\ref{sec:kg-results}. Nama di dalam
tanda petik balik disisipkan menurut tipe \emph{resource} dan menurut medannya.

\begin{lstlisting}[caption={Jalur tulis \emph{loader} (\texttt{loader.py}).},label={lst:loader-merge}]
// nodes -- one statement per FHIR resource type
UNWIND $rows AS row
MERGE (n:Resource {uid: row.uid})
SET n += row.props SET n:`<ResourceType>`

// shared terminology nodes, keyed on (system, code)
UNWIND $rows AS row
MERGE (c:Code {system: row.system, code: row.code})
ON CREATE SET c.display = row.display

// resource-to-resource references
UNWIND $rows AS row
MATCH (s:Resource {uid: row.src})
MATCH (t:Resource {uid: row.dst})
MERGE (s)-[:`<FIELD>`]->(t)
RETURN count(*) AS n      // a shortfall means a MATCH found no node

// resource-to-code edges
UNWIND $rows AS row
MATCH (s:Resource {uid: row.src})
MATCH (c:Code {system: row.system, code: row.code})
MERGE (s)-[:`<FIELD>`]->(c)
\end{lstlisting}

\subsection{Predikat \emph{Observation Gate}}
Tiap entri \texttt{OBS\_GATE} merupakan predikat atas fitur turunan graf. \texttt{decide}
membiarkan predikat tersebut \emph{menimpa} keputusan gerbang-dan-\texttt{AND} pada graf,
bukan mengirisnya, sehingga cabang disjungtif seperti \texttt{RAAI:1} justru menambah
pasien alih-alih menyaringnya keluar.

\begin{lstlisting}[language=Python,caption={\emph{Observation gate} (\texttt{obs\_gate.py}), diringkas.},label={lst:obs-gate}]
OBS_GATE = {
    "CCD-2023:BB:1":    lambda f: f["lvef_min"] is not None and f["lvef_min"] <= 40,
    "CCD-2023:BB:2":    lambda f: f["lvef_min"] is not None and f["lvef_min"] < 50,
    "CCD-2023:SGLT2:4": lambda f: f["hf"] and f["lvef_min"] is not None
                                  and f["lvef_min"] <= 40,
    "CCD-2023:RAAI:1":  lambda f: f["htn"] or f["dm"] or f["ckd"]
                                  or (f["lvef_min"] is not None and f["lvef_min"] <= 40),
    "CCD-2023:WM:2":    lambda f: f["bmi_ge25"],
    # ... ten further predicates, one per observation-gated recommendation
}
OBS_GATED = set(OBS_GATE)

def decide(graph_applicable, feats):
    keep = {r for r in graph_applicable if r not in OBS_GATED}
    if feats.get("ccd"):          # obs-gated recs fire only for CCD patients
        keep |= {r for r, pred in OBS_GATE.items() if pred(feats)}
    return keep
\end{lstlisting}

\subsection{\emph{Prompt} \emph{Briefing}}
\emph{Prompt} inilah yang mengikat model bahasa pada fakta hasil penelusuran; pemeriksaan
\emph{faithfulness} pada Subbab~\ref{sec:faithfulness} kemudian memverifikasi bahwa ikatan
tersebut bertahan.

\begin{lstlisting}[caption={\emph{Prompt} \emph{grounded briefing} (\texttt{crs\_briefing.py}), diringkas.},label={lst:briefing-prompt}]
SYSTEM
You are summarizing decision-support facts retrieved from a clinical knowledge
graph for a physician treating a patient with chronic coronary disease (CCD).
The briefing INFORMS the clinician; it does not make the decision.

USER
Write a clinician-facing briefing from the JSON facts below, in TWO clearly-headed
parts: a MAIN BRIEFING a cardiologist reads in 30 seconds, then a DETAILS chapter
for drill-down.

CONTENT RULES (strict -- these preserve trust):
- Use ONLY the facts provided. Never add recommendations, drugs, doses, or
  evidence that are not present in the JSON.
- Justify applicability ONLY from `requiresAlso` and the patient's actual
  `conditions` / `medications`. A rec's OTHER text preconditions are NOT
  established facts about this patient. If a rec's action assumes such an
  unconfirmed precondition, phrase the action CONDITIONALLY, never as an
  assertion.
- Every item in `applicable` and `safetyFlags` MUST appear at least once, named
  by its recommendation id. You may GROUP and COMPRESS them, but never drop one.
- Never alter a COR value, and never attach a PMID to the wrong recommendation.
- State a count only if it equals the number of items given.
- If a list is empty, say so explicitly -- do not fill the gap.

[... the two part templates follow, then the JSON facts payload ...]
\end{lstlisting}


\section{Contoh \emph{Grounded Briefing}}
\label{app:briefing}

Listing~\ref{lst:briefing-example} memuat satu \emph{grounded briefing} secara utuh, persis
sebagaimana dihasilkan sistem dan disimpan sebagai
\texttt{eval/briefing\_22f599b6-\ldots.md}. Pasiennya adalah rekam medis sintetis yang sama
dengan yang dipakai sebagai contoh \emph{resource} FHIR pada Bab~II. Tidak ada bagian yang
disunting, diurutkan ulang, maupun dipendekkan: inilah keluaran yang dinilai pada
Subbab~\ref{sec:faithfulness}.

\begin{lstlisting}[basicstyle=\scriptsize\ttfamily,numbers=none,caption={\emph{Grounded briefing} lengkap bagi satu pasien, sebagaimana dibangkitkan (\texttt{claude-sonnet-5}) dan disimpan pada \texttt{eval/}.},label={lst:briefing-example}]
# MAIN BRIEFING

**1. BOTTOM LINE**
This is a CCD patient with established atherosclerotic disease (ischemic heart disease, history of STEMI/MI, history of CABG), hypertension, diabetic kidney disease, and obesity. Top priorities: refer to cardiac rehabilitation (COR 1) and treat hypertension to a <130/<80 target with GDMT (COR 1); antiplatelet strategy should reflect PCI history but any oral-anticoagulant-dependent regimen must first be confirmed clinically. Key safety concerns: nitrate/PDE5-inhibitor interaction (harm) and prasugrel contraindication if stroke/TIA/ICH history, plus two no-benefit flags on beta-blocker and clopidogrel use.

**2. SAFETY FLAGS (stop-and-check)**
- Metoprolol succinate ER / metoprolol tartrate -- no MACE benefit if no prior MI or LVEF<=50% and no other BB indication [BB:4, 3-NoBenefit]
- Clopidogrel -- not useful added to aspirin without recent ACS or PCI-related DAPT indication [APT:7, 3-NoBenefit]
- Nitroglycerin (spray and sublingual) -- do not combine with PDE5 inhibitors: risk of severe hypotension [SHA:3, 3-Harm]
- Prasugrel -- contraindicated with history of stroke/TIA/ICH [APT:9, 3-Harm]

**3. ACTIONS FOR THIS PATIENT**
- Refer to cardiac rehabilitation program [CR:1:1, COR 1]
- Hypertension lifestyle education toward BP <130/<80 [BPM:2, COR 1]
- First-line antihypertensive GDMT -- ACE inhibitor, ARB, or beta blocker (e.g., lisinopril, amlodipine, metoprolol succinate) (+17 more options), add-on agents as needed [BPM:3, COR 1]
- If on/requiring oral anticoagulation: DAPT 1-4 weeks then clopidogrel alone for 6 months, added to DOAC [APT:11, COR 1]
- Weight management counseling on diet, lifestyle, and weight-loss goals [WM:2, COR 1]
- If on/requiring oral anticoagulation and high thrombotic/low bleeding risk: continue aspirin + clopidogrel up to 1 month [APT:12, COR 2a]
- Consider GLP-1 receptor agonist (semaglutide preferred) for further weight reduction, with diet/activity counseling [WM:3, COR 2a]

---

# DETAILS & RATIONALE

**1. FURTHER PATIENT-SPECIFIC OPTIONS (COR 2b -- "may be considered")**
- Extended DAPT (aspirin + clopidogrel) beyond 12 months up to 3 years [APT:4, COR 2b]
- Vorapaxar added to aspirin (no stroke/TIA/ICH history) [APT:5, COR 2b]
- DAPT after CABG to reduce saphenous vein graft occlusion [APT:6, COR 2b]
- If on/requiring oral anticoagulation with low atherothrombotic risk: discontinue aspirin, continue DOAC alone (rivaroxaban, apixaban, dabigatran, edoxaban) at 1 year post-PCI [APT:13, COR 2b]
- Reassess ongoing need for long-term (>1 year) beta-blocker therapy [BB:3, COR 2b]

**2. WHY THESE APPLY**
- CR:1:1 -- triggered by History of MI, History of CABG, History of PCI
- BPM:2 -- triggered by Essential Hypertension
- BPM:3 -- triggered by Essential Hypertension
- APT:11 -- triggered by History of PCI (anticoagulant need not confirmed in patient's medication list)
- WM:2 -- triggered by BMI 30+/obesity finding
- APT:12 -- triggered by History of PCI (anticoagulant need not confirmed)
- WM:3 -- triggered by BMI 30+/obesity finding
- APT:4 -- triggered by History of MI
- APT:5 -- triggered by History of MI
- APT:6 -- triggered by History of CABG
- APT:13 -- triggered by History of PCI (anticoagulant need not confirmed)
- BB:3 -- triggered by History of MI

**3. STANDARD CCD CARE (applies to all CCD patients)**
- Nutrition: heart-healthy diet [NUTRI:1, COR 1]; reduced saturated fat [NUTRI:2, COR 2a]; low sodium/processed meat limits [NUTRI:3, COR 2a]; limit refined carbs/sugar drinks [NUTRI:4, COR 2a]
- Tobacco: assess use at every visit [TP:1, COR 1]
- Alcohol/Substance use: routine counseling [ASU:1, COR 1]; limit alcohol intake [ASU:2, COR 2a]
- Lipid management: high-intensity statin [LIPID:1, COR 1]; moderate-intensity statin if high-intensity not tolerated [LIPID:2, COR 1]; periodic lipid panel monitoring [LIPID:3, COR 1]; consider ezetimibe [LIPID:10, COR 2b]; consider bempedoic acid/inclisiran [LIPID:11, COR 2b]
- Weight management: routine BMI/waist circumference assessment [WM:1, COR 1]
- Cardiac rehabilitation: also indicated after SCAD event (not this patient's trigger) [CR:1:3, COR 1]
- Physical activity: aerobic exercise >=150 min/wk [PA:1, COR 1]; resistance training >=2 days/wk [PA:2, COR 1]; reduce sedentary time [PA:3, COR 2a]
- Antiplatelet therapy: low-dose aspirin if no anticoagulant indication [APT:1, COR 1]; DAPT 6 months post-PCI then SAPT [APT:2, COR 1]; P2Y12 monotherapy after short DAPT with DES [APT:3, COR 2a]; low-dose rivaroxaban + aspirin in select high-risk patients [APT:15, COR 2a]; PPI with DAPT for GI bleeding risk [APT:16, COR 2a]; DOAC monotherapy if no acute antiplatelet indication [APT:14, COR 2b]
- Beta blockers: recommended if LVEF<=40% [BB:1, COR 1]; preferred agents (metoprolol succinate, carvedilol, bisoprolol) if LVEF<50% [BB:2, COR 1]
- Renin-Angiotensin-Aldosterone inhibitors: ACEi/ARB for hypertension, diabetes, LVEF<=40%, or CKD [RAAI:1, COR 1]
- Immunizations: annual influenza vaccine [IMMU:1, COR 1]; COVID-19 vaccination [IMMU:2, COR 1]; pneumococcal vaccine [IMMU:3, COR 2a]
- Angina therapy: beta blocker/CCB/long-acting nitrate first-line [MTRA:1, COR 1]; add second class if symptomatic [MTRA:2, COR 1]; ranolazine if still symptomatic [MTRA:3, COR 1]; sublingual/spray nitroglycerin for immediate relief [MTRA:4, COR 1]
- Mental health: screening [MHC:1, COR 2a]; treatment if indicated [MHC:2, COR 2a]
- Sexual health: individualize resumption of activity [SHA:1, COR 2a]; cardiac rehab/exercise to reduce risk [SHA:2, COR 2a]
- Environmental exposures: minimize air pollution exposure [EE:1, COR 2a]; minimize climate-related exposures [EE:2, COR 2b]
- Colchicine: may be considered for secondary prevention [COLC:1, COR 2b]
- Refractory angina: enhanced external counterpulsation if no other options [MRA:1, COR 2b]

**4. LITERATURE-SUPPORTED OPTIONS**
- Disorder of kidney due to diabetes mellitus + canagliflozin 100 MG: pmi 1.249, 4 articles
- Metabolic syndrome X + Simvastatin (10/20/40 MG): pmi 0.759, 6 articles (each strength)
- Myocardial infarction + irbesartan (75/150 MG): pmi 0.537, 4 articles
- Myocardial infarction + ramipril (2.5/5/10 MG): pmi 0.475, 39 articles
- Myocardial infarction + ticagrelor 90 MG: pmi 0.43, 68 articles
- Myocardial infarction + prasugrel 10 MG: pmi 0.383, 47 articles
- Myocardial infarction + clopidogrel (75/300 MG): pmi 0.367, 148 articles
- Myocardial infarction + aspirin 81 MG (tablet/capsule): pmi 0.353, 249 articles

**5. EVIDENCE (recId -- citedPmids)**
- CR:1:1 -- 17387127, 32100573, 32227755, 34355130, 34587751
- BPM:2 -- 26551272, 28564682, 30919899, 31169350, 32866419
- BPM:3 -- 10639539, 15351192, 19454737, 26551272, 9796841
- APT:11 -- 32101251
- APT:12 -- 32101251, 31707833
- APT:4 -- 25773268, 26324537
- APT:5 -- 22443427, 22932716, 26338971
- APT:6 -- 30072154
- APT:13 -- 31475793
- BB:3 -- 24927909, 25511558, 28571635, 32542362, 33428707, 34227397, 33877862
- WM:2 -- 18056904, 19744231, 25369490, 30326501
- WM:3 -- 26132939, 30122305, 33567185, 35015037
- NUTRI:1 -- 23999993, 30864165, 34255848, 9989963
- NUTRI:2 -- 23999993, 24355537, 30864165, 32428300, 34255848, 9989963
- NUTRI:3 -- 22110105, 23999993, 32428300, 9989963
- NUTRI:4 -- 11136953, 23999993, 32428300, 34255848, 9989963
- TP:1 -- 18617085, 22695462, 28056007
- ASU:1 -- 25986138, 28793133, 29535062, 31976871, 21156654
- ASU:2 -- 20338495, 22076059, 35333364
- LIPID:1 -- 15755765, 21067804, 30712900
- LIPID:2 -- 12114037, 21067804, 7968073, 8801446, 8992351, 9841303
- LIPID:3 -- 21067804, 27533947, 31952841, 34861203
- LIPID:10 -- 26039521, 28231942, 28768335, 28465286
- LIPID:11 -- 30865796, 31714986, 32187462
- WM:1 -- 10511607, 16820550, 17846390, 20038741, 20696950, 29490333
- CR:1:3 -- 17387127, 26923234, 27055757, 29472380, 32227755, 33884107, 34355130
- PA:1 -- 15007010, 25488550, 29654095
- PA:2 -- 21450617, 21896934, 34805309
- PA:3 -- 32563261, 32843054, 33930434
- APT:1 -- 11786451, 19482214, 33999548
- APT:2 -- 22179532, 25236346, 25616646, 25777667
- APT:3 -- 22999717, 24177257, 28110296, 32551860, 32795096
- APT:14 -- 27959713, 31603196, 33250267
- APT:15 -- 28844192, 32223318, 32436455
- APT:16 -- 34408652
- BB:1 -- 11356434, 12390947, 32289158
- BB:2 -- 10376614, 11356434, 11386264, 11835035, 12390947, 12853193, 7600653
- RAAI:1 -- 10821360, 1386652, 1463530, 7477219, 8104270
- IMMU:1 -- 24150467, 25940444, 26310262, 31444266, 32535214, 33275134, 33719496
- IMMU:2 -- 33622834, 35076665, 33587998
- IMMU:3 -- 22739637, 32259237, 34521082
- MTRA:1 -- 10349897, 1672481, 2178381
- MTRA:2 -- 10349897, 24776376, 30165445, 8557899
- MTRA:3 -- 14734593, 16875985
- MTRA:4 -- 1577055, 27468079
- MHC:1 -- 23759474, 31633746, 31907104, 30698764
- MHC:2 -- 28209727, 30043065, 31907104, 30698764
- SHA:1 -- 17628195, 21427375
- SHA:2 -- 8618365
- EE:1 -- 16893917, 25810496, 26613345, 27233746, 31433918, 31492503, 32384243
- EE:2 -- 31376629, 33301532, 34098271
- COLC:1 -- 23265346, 32865380
- MRA:1 -- 10362181
- Safety flags (BB:4, APT:7, SHA:3, APT:9) -- no citedPmids provided in source data
\end{lstlisting}
